
% 开始写正文
% 设置正文的页边距
\newgeometry{top=2cm, left=2.5cm, right=2.5cm}
% 设置正文的页眉页脚
\fancyhf{}
\fancyhead[C]{ }
% 此处修改右上角页码
%\fancyhead[R]{Page \thepage\ of 18}
\rhead{Page \thepage{} of \pageref{LastPage}}
%\fancyhead[R]{Page \thepage\ of 15}
\fancyhead[L]{Team \# 1917694}
\fancyfoot[C]{\bfseries\thepage}


\section{Introduction}%第一节 引言（机器自动编号） 

{\bf Theorem 1.} {\it For any fixed $\rho>0$, we have 
	
	$$||R(u,\rho)||\leq ||R_1(u,\rho)||\leq 2||R(u,\rho)||,
	~~~\forall u\in H. 
	\eqno(10)
	$$ 
	Hence, (2) and (3) are equivalent for any fixed$
	\rho>0$.}	 %一般定理的写法
\subsection{Problem Background}

\subsection{Restatement of the problem}
Describe the origin of the problem according to your own understanding.If the topic has biology or other things, you can add pictures or other things here.
\begin{itemize}
	\item Build a mathematical model to 
	\item based upon..., use your model to
	\item Inlight of your predictive analysis,...
	\item use your model to
\end{itemize}
\cite{__2020}
\cite{_pvar_2020}

\section{Model Assumptions and Symbols}

\subsection{Assumptions and Justifications}

To simplify the problem , we make the following basic assumptions.

\begin{itemize}		%无序列表
	\item Assumption1
	\item Assumption2
\end{itemize}

\begin{enumerate}	%有序列表
	\item Assumption1
	\item Assumption2
\end{enumerate}

\subsection{Symbols}
Table 1 shows symbols and notaions used in this paper.

\begin{center}	%表格，千万不要自己写，累死人
	\begin{tabular}{lc}
		\hline
		Symbol & definition \\ \hline
		a&the symbol of            \\
		b&the definition of            \\
		c&            \\
		d&            \\
		e&            \\
		f&            \\
		g&            \\
		h&            \\ \hline
	\end{tabular}
\end{center}

\section{Model Overview}


\begin{figure}
	\centering
	\includegraphics[scale=0.5]{test1.eps}\\%插入图片，别忘了要在上面导入包
	\caption{李佳迅}	%图注
\end{figure}


\section{The Data}
\subsection{Data Collection}
this is the description of data collection
\subsection{Data Processing}
this is the description data processing

\section{Model I}
\subsection{Description of ...}
this is the description of model
对题目中物理量的描述
\subsection{...Model}
this is the description of data collection
模型I中用到的模型方法，简单叙述一下模型方法，将模型方法的概念和文中的物理量联系在一起，公式一定要居中

\begin{equation}
	\begin{split}
		\frac{1}{2}=0.5
	\end{split}
\end{equation}




\subsection{Result}
this is the result of our work,我们在什么什么的基础到我会的哦阿尼法黑of活期为我偶尔服你哦亲问佛我hi二分为hi  u让我去哈佛我还欠我一饿哦哈佛和后悔维护费熊前后好熊顽强哈诶of或氰化物ifhi
\section{Model II}


\section{Model III}


\section{Test the Model}


\section{Conclusion}

\subsection{summary of results}

\subsubsection{result of problem1}

\subsubsection{result of problem2}


\subsection{Strength}
\begin{itemize}
	\item ........
	\item ........
	\item ........
\end{itemize}
\subsection{possible Improvements}
不宜太多



